%% Part IV: Inside Tokio. Run queues, driver duty, mio, and the second kind of wake.

\subsection{The runtime at a glance}
\label{sec:wakers-ch10}

\textit{Tokio's default runtime spreads tasks across several worker threads, each with its own queue, and balances the load by stealing. Figure~\ref{fig:wakers-10.1} shows the whole structure.}

\begin{figure}[htbp]
    \centering
    % Coordinates follow a 940 x 660 grid; y grows downwards.
    \newcommand{\wkworker}[6]{% x, name, queued tasks (0-8), pill colour, pill text, notes
        \draw[rounded corners=6pt, draw=gray!50, fill=white, line width=0.6pt] (#1,152) rectangle ++(280,200);
        \node[anchor=north west, inner sep=0] at ({#1+16},162) {\small\textbf{#2}\quad\textcolor{gray}{one OS thread}};
        \fill[rounded corners=3pt, #4!15] ({#1+16},186) rectangle ++(248,26);
        \node[anchor=west, font=\ttfamily\scriptsize, text=#4!70!black] at ({#1+24},199) {#5};
        \node[anchor=west, inner sep=0, text=gray] at ({#1+16},230) {LIFO slot};
        \node[anchor=west, inner sep=0, text=gray] at ({#1+76},230) {local run queue (up to 256)};
        \ifnum#3>0 \fill[rounded corners=1pt, StmtProposition] ({#1+16},244) rectangle ++(22,22);
        \else \draw[rounded corners=1pt, gray!50] ({#1+16},244) rectangle ++(22,22); \fi
        \foreach \i in {0,...,7}{
            \ifnum\i<#3 \fill[rounded corners=1pt, StmtProposition] ({#1+76+20*\i},247) rectangle ++(16,16);
            \else \draw[rounded corners=1pt, gray!50] ({#1+76+20*\i},247) rectangle ++(16,16); \fi
        }
        \node[anchor=west, text=gray] at ({#1+236},255) {\dots};
        \node[anchor=north west, inner sep=0, text width=4.3cm, text=gray] at ({#1+16},280) {#6};
    }
    \figurefonts\resizebox{\linewidth}{!}{\begin{tikzpicture}[x=0.017cm, y=-0.019cm, font=\sffamily\scriptsize, >={Stealth}]
        % other threads
        \draw[rounded corners=6pt, draw=StmtAxiom, fill=StmtAxiom!8, line width=0.6pt] (20,20) rectangle (290,106);
        \node[anchor=north west, inner sep=0, text width=4.3cm] at (36,32) {{\small\textbf{Other threads}}\\[2pt] \textcolor{gray}{blocking pool, main(), any non-worker}\\[2pt] \texttt{spawn()} $\cdot$ \texttt{wake()}};
        \draw[->, gray, thick] (290,63) -- (326,63);
        % global inject queue
        \draw[rounded corners=6pt, draw=StmtProposition, fill=orange!8, line width=0.6pt] (330,20) rectangle (920,106);
        \node[anchor=north west, inner sep=0] at (348,32) {\small\textbf{Global inject queue}\quad\textcolor{gray}{\scriptsize shared by all workers, behind a mutex}};
        \foreach \i in {0,1,2}{\fill[rounded corners=1pt, StmtProposition] ({348+24*\i},64) rectangle ++(18,18);}
        \foreach \i in {3,4,5}{\draw[rounded corners=1pt, gray!50] ({348+24*\i},64) rectangle ++(18,18);}
        \node[anchor=north west, inner sep=0, text=gray, text width=6.8cm] at (510,60) {in: remote wakes, outside spawns, overflow\\ out: workers pull when idle, and periodically};
        % inject queue to workers
        \draw[->, gray, thick] (400,106) .. controls (400,130) and (160,124) .. (160,150);
        \draw[->, gray, thick] (470,106) -- (470,150);
        \draw[->, gray, thick] (560,106) .. controls (560,130) and (780,124) .. (780,150);
        % workers
        \wkworker{20}{Worker 0}{6}{Primary1}{executor $\cdot$ polling task}{next task: LIFO slot first (max 3 in a row), then the local queue.\\[2pt] Full? Half moves to the inject queue.}
        \wkworker{330}{Worker 1}{4}{Primary1}{executor $\cdot$ polling task}{A task woken from this thread lands in this worker's LIFO slot.\\[2pt] (LIFO-slot tasks can't be stolen.)}
        \wkworker{640}{Worker 2}{0}{StmtProposition}{idle $\cdot$ running the driver}{Empty: 1) check inject queue, 2) steal from a sibling, 3) park.\\[2pt] (First to park takes the driver.)}
        % stealing
        \draw[->, StmtProposition, thick, dashed] (560,352) .. controls (580,392) and (690,392) .. (712,356);
        \node[text=StmtProposition, font=\sffamily\scriptsize\bfseries] at (600,402) {steal half of the local queue};
        % driver
        \draw[rounded corners=6pt, draw=orange!70!black, fill=yellow!10, line width=0.6pt] (330,420) rectangle (920,570);
        \node[anchor=north west, inner sep=0] at (348,430) {\small\textbf{Driver}\quad\textcolor{gray}{\scriptsize one worker holds it at a time; other idle workers sleep on a condvar}};
        \draw[rounded corners=4pt, draw=gray!50, fill=white] (348,456) rectangle (618,530);
        \node[anchor=north west, inner sep=0, text width=4.3cm] at (360,466) {\texttt{I/O driver (mio)}\\[2pt] \textcolor{gray}{epoll set $\rightarrow$ per-socket ScheduledIo; reader and writer wakers live here}};
        \draw[rounded corners=4pt, draw=gray!50, fill=white] (632,456) rectangle (902,530);
        \node[anchor=north west, inner sep=0, text width=4.3cm] at (644,466) {\texttt{Time driver}\\[2pt] \textcolor{gray}{timer wheel, 6 levels $\times$ 64 slots; sets the epoll\_wait timeout}};
        \node[anchor=north west, inner sep=0, text=gray, text width=9.6cm] at (348,542) {Busy workers also check it without blocking every 61 ticks, so I/O isn't starved.};
        \draw[->, gray, thick] (760,352) -- (760,418) node[midway, right, text=gray] {parks};
        \draw[->, StmtLemma, thick] (860,418) -- (860,354) node[midway, right, text=StmtLemma, font=\sffamily\scriptsize\bfseries] {wake()};
        % kernel
        \draw[rounded corners=6pt, draw=gray!50, fill=gray!8, line width=0.6pt] (330,600) rectangle (920,648);
        \node[anchor=west] at (340,624) {\small\textbf{Kernel}\quad\texttt{\scriptsize epoll\_wait $\cdot$ sockets $\cdot$ eventfd (unpark signal)}};
        \draw[<->, gray, thick] (560,572) -- (560,598);
        % owned tasks
        \draw[rounded corners=6pt, draw=gray!60, dashed, fill=white, line width=0.6pt] (20,420) rectangle (300,648);
        \node[anchor=north west, inner sep=0, text width=4.3cm] at (36,432) {\small\textbf{OwnedTasks}\\[4pt] \scriptsize\textcolor{gray}{A list of every live task, including idle ones parked on a waker.\\[4pt] Not a queue. Never scanned to decide what to poll; used only for abort, shutdown and metrics.\\[4pt] A task that's not queued and not woken just sits here.}};
    \end{tikzpicture}}
    \caption{Tokio's multi-thread runtime, simplified. The \texttt{current\_thread} runtime of Section~\ref{sec:wakers-ch5} is the same picture with one worker and no stealing.}
    \label{fig:wakers-10.1}
\end{figure}

The parts of the figure:

\begin{description}
    \item[Worker] One operating-system thread that polls tasks, and goes on driver duty when it has none.
    \item[Local run queue] Each worker's own queue of up to 256 runnable tasks.
    \item[LIFO slot] Each worker's single-task ``next up'' slot, filled by local wakes.
    \item[Inject queue] A queue shared by all workers. It receives remote wakes, tasks spawned from outside the workers, and overflow from full local queues.
    \item[Driver] The I/O and timer machinery, run by one idle worker at a time (Section~\ref{sec:wakers-ch11}).
    \item[Blocking pool] Up to 512 extra threads by default, for \code{spawn\_blocking} and file I/O. They never poll async tasks; they only call wakers.
    \item[List of all tasks] Every live task, idle ones included. It is not a queue and is never searched to decide what to run; it exists for cancellation, shutdown and metrics. A stranded task sits here.
\end{description}

Where a woken task lands depends on who woke it:

\begin{table}[htbp]
    \centering\small
    \begin{tabularx}{\linewidth}{>{\raggedright\arraybackslash}p{2.4cm} Y >{\raggedright\arraybackslash}p{2.6cm}}
        \toprule
        \textbf{Kind of wake} & \textbf{Where the task goes} & \textbf{In Section~\ref{sec:wakers-ch5}} \\
        \midrule
        Local wake & That worker's LIFO slot. The slot's previous occupant moves to the back of the local queue. A full local queue sends half its tasks to the inject queue. & Task A \\
        Remote wake & The inject queue. Then a sleeping worker is woken: through its condition variable, or through the eventfd if it is on driver duty (Lemma~\ref{lem:remote-wakes-need-a-poke}). & Task B \\
        \bottomrule
    \end{tabularx}
    \caption{Where a woken task lands.}
    \label{tab:wakers-10.1}
\end{table}

Each structure exists for a reason:

\begin{itemize}
    \item \textbf{The LIFO slot}\index{LIFO slot} exploits a common pattern. When task X sends a message that wakes task Y, Y usually needs data X has just touched, so running Y next on the same core finds that data still in cache.
    \begin{itemize}
        \item At most three LIFO polls run in a row, so two tasks can't ping-pong forever.
        \item A task in the LIFO slot cannot be stolen.
    \end{itemize}
    \item \textbf{Work stealing}\index{Work stealing} handles uneven load. A worker whose own queue and the inject queue are both empty takes half of a busy sibling's local queue before sleeping.
    \item \textbf{Task selection order:} the LIFO slot, then the local queue. Every so often the inject queue goes first, for fairness.
    \item \textbf{I/O fairness:} a busy worker checks the driver without blocking every 61 tasks, so a full queue cannot starve I/O.
\end{itemize}

\begin{quotebox}{Summary}
    \begin{itemize}
        \item Each worker has a LIFO slot and a 256-task local queue; a shared inject queue takes remote wakes.
        \item Idle workers steal half of a sibling's queue. One idle worker at a time is on driver duty.
        \item The list of all tasks is never used for scheduling.
    \end{itemize}
\end{quotebox}

\subsection{Driver duty}
\label{sec:wakers-ch11}

\textit{``Reactor'' is the traditional name for the part of a runtime that waits for events. In Tokio it isn't a component or a thread. It is a duty that a worker takes up when it has nothing else to do.}

\begin{definition}{Driver}{driver}
    \index{Driver}\index{Reactor}
    Tokio's \emph{driver} is the private machinery that waits on the operating system and dispatches what it reports: socket readiness, expired timers, signals and child-process exits. It has no thread of its own; whichever worker goes on driver duty (Definition~\ref{def:driver-duty}) runs it.
\end{definition}

A little history clarifies the names:

\begin{itemize}
    \item The \emph{reactor pattern} is old: block on an operating-system event queue, then dispatch each event to whoever registered interest.
    \item Before version 1.0, Tokio exposed it as \code{tokio\_core::reactor::Core}.
    \item Modern Tokio has no public reactor type and no reactor thread. Its documentation calls the same machinery the \emph{driver}.
    \item In the multi-thread runtime, the first idle worker to take the driver's lock goes on driver duty; other idle workers sleep on a condition variable.
\end{itemize}

\subsubsection{Two different things called ``poll''}

Tokio sits on a library called mio (Section~\ref{sec:wakers-ch12}), whose main type is \code{Poll}, with a method also named \code{poll}. The two have nothing in common:

\begin{table}[htbp]
    \centering\small
    \begin{tabularx}{\linewidth}{>{\raggedright\arraybackslash}p{1.8cm} >{\raggedright\arraybackslash}p{4.6cm} Y}
        \toprule
         & \textbf{\texttt{Future::poll(cx)}} & \textbf{\texttt{mio::Poll::poll(\&mut events, timeout)}} \\
        \midrule
        Meaning & ``Make progress if you can'' & A thin wrapper around \texttt{epoll\_wait} \\
        Blocks? & Never; returns \texttt{Pending} instead & Yes, until an event or the timeout \\
        Called by & The executor, on a task & The driver, in exactly one place \\
        When & When the task is on the run queue & \textbf{Blocking} on driver duty (timeout: next timer, or none). \textbf{Non-blocking} (zero timeout) every 61 tasks while busy. Never in a busy loop. \\
        \bottomrule
    \end{tabularx}
    \caption{The two meanings of ``poll.''}
    \label{tab:wakers-11.1}
\end{table}

\begin{lemma}{Single blocking point}{single-blocking-point}
    A Tokio runtime waits for I/O, timers, signals and child processes in exactly one place: the call to \code{mio::Poll::poll} made by the worker on driver duty. (If the runtime was built without I/O, it is a plain thread park with a timeout instead.) Other idle workers wait on a condition variable for work, not for events.
    \why The four drivers below are layered so that each wraps the next. A park request passes down through all of them, and only the innermost blocks.
\end{lemma}

\begin{lstlisting}[language={}, caption={The four driver layers, from outermost to innermost.}, label=lst:wakers-11.1]
idle worker -> park()
  -> Time driver      next deadline is 37 ms away -> park_timeout(37 ms);
                      afterwards, fire every expired timer
    -> Process driver   (Unix) reap exited child processes whose handle was dropped
      -> Signal driver    (Unix) a signal handler wrote a byte to a registered socket
        -> I/O driver       mio::Poll::poll(events, 37 ms) -> epoll_wait(..., 37)   <- blocks here
                            for each event: token -> that socket's state -> wake()
\end{lstlisting}

\begin{itemize}
    \item \textbf{I/O driver:} any descriptor epoll can watch; network sockets are only one kind.
    \item \textbf{Signal and process drivers:} thin layers that register extra descriptors with the I/O driver.
    \item \textbf{Time driver:} turns the timer wheel's next deadline into the timeout (Section~\ref{sec:wakers-ch15}).
    \item \textbf{On Windows:} only the I/O driver (IOCP) and the time driver exist.
    \item \textbf{Files} never pass through any driver (Axiom~\ref{ax:files-always-ready}).
\end{itemize}

\subsubsection{The single-threaded runtime}

The \code{current\_thread} runtime has no worker threads. Its only executor is whichever thread calls \code{block\_on}, typically the main thread, and that thread owns the driver outright, with no lock.

\begin{lstlisting}[language={}, caption={The \texttt{current\_thread} scheduler loop, paraphrased from Tokio's source.}, label=lst:wakers-11.2]
loop {
    if the main future was woken: poll it      // the block_on future is not a spawned task
    repeat up to 61 times:
        take the next task (local queue; the inject queue every 31 tasks)
        nothing to run?  -> driver duty: epoll_wait(timeout = next timer); restart loop
        otherwise        -> poll that task
    ran 61 tasks without parking -> check the driver with a zero timeout
}
\end{lstlisting}

\begin{table}[htbp]
    \centering\small
    \begin{tabularx}{\linewidth}{>{\raggedright\arraybackslash}p{4.4cm} Y Y}
        \toprule
         & \textbf{Multi-thread} & \textbf{\texttt{current\_thread}} \\
        \midrule
        Who polls tasks & N worker threads & The thread inside \texttt{block\_on} \\
        Who has driver duty & Whichever idle worker takes the lock & That same thread, with no lock \\
        A remote wake & Inject queue; a worker is unparked & Inject queue; the eventfd is written \\
        Progress when nobody is in \texttt{block\_on} & Yes & \textbf{None} \\
        \bottomrule
    \end{tabularx}
    \caption{The two runtime flavours compared.}
    \label{tab:wakers-11.2}
\end{table}

\begin{corollary}{No block\_on, no progress}{no-block-on-no-progress}
    A \code{current\_thread} runtime polls no tasks, handles no I/O and fires no timers while no thread is inside its \code{block\_on}.
\end{corollary}

In practice:

\begin{itemize}
    \item Tasks spawned onto such a runtime through a handle from another thread will not run until some thread enters \code{block\_on}.
    \item If two threads call \code{block\_on} on the same \code{current\_thread} runtime, only one holds the scheduler at a time; the other waits.
\end{itemize}

\begin{quotebox}{Summary}
    \begin{itemize}
        \item The ``reactor'' is the driver, run as a duty by an idle worker (Definition~\ref{def:driver}).
        \item \code{Future::poll} never blocks; mio's \code{Poll::poll} is the single blocking point (Lemma~\ref{lem:single-blocking-point}).
        \item A \code{current\_thread} runtime makes progress only inside \code{block\_on} (Corollary~\ref{cor:no-block-on-no-progress}).
    \end{itemize}
\end{quotebox}

\subsection{mio: where Tokio meets the operating system}
\label{sec:wakers-ch12}

\textit{Between Tokio and the kernel sits a small library that knows nothing about futures. Seeing exactly where it begins and ends removes most of the mystery from async I/O.}

The work divides into two questions:

\begin{itemize}
    \item \emph{Which handles are ready?} That is mio's job. mio\index{mio} is a thin, synchronous Rust wrapper over each operating system's readiness call.
    \item \emph{Which task is waiting on that handle, and how do I wake it?} That is Tokio's job.
\end{itemize}

\begin{definition}{Token}{token}
    \index{Token (mio)}
    A \emph{token} is a number (\code{Token(usize)}) that Tokio attaches to a descriptor when it registers it with mio. Tokio uses a pointer to the socket's internal state, including its waker slots, as the token.
\end{definition}

\begin{lemma}{Token round trip}{token-round-trip}
    The token passes through the kernel unchanged: it is stored with the registration and returned with every event for that descriptor. The driver therefore goes from an event to the right waker slot with no lookup.
    \why mio places the token in the \code{data} field of the epoll registration, which the kernel returns verbatim. Listing~\ref{lst:wakers-12.1} shows it happening.
\end{lemma}

\begin{figure}[htbp]
    \centering\small
    \begin{tabularx}{\linewidth}{>{\raggedright\arraybackslash}p{2.6cm} Y}
        \toprule
        \textbf{your code} & \texttt{sock.read(\&mut buf).await} \\
        \midrule
        \textcolor{Primary3}{\textbf{Tokio}} \newline {\footnotesize\color{gray} futures, tasks, wakers} &
            \texttt{TcpStream::poll\_read(cx)} checks readiness, otherwise stores \texttt{cx.waker()}. \newline
            \texttt{ScheduledIo}: per socket, readiness flags plus reader and writer wakers. \newline
            Driver: for each event, \textbf{token $\rightarrow$ ScheduledIo $\rightarrow$ \texttt{wake()}}. \\
        \midrule
        \textcolor{StmtProposition}{\textbf{the handoff}} &
            \textbf{\texttt{Token(usize)}}, chosen by Tokio when it registers a socket and returned unchanged by mio. Tokio uses a pointer to the socket's \texttt{ScheduledIo}, so no lookup is needed. \\
        \midrule
        \textcolor{StmtLemma}{\textbf{mio}} \newline {\footnotesize\color{gray} no futures, no wakers} &
            \texttt{Registry::register(\&mut sock, Token(n), Interest::READABLE)} \newline
            \texttt{Poll::poll(\&mut events, timeout)} $\rightarrow$ \texttt{[(Token, readable / writable), ...]} \newline
            Non-blocking socket types $\cdot$ \texttt{mio::Waker} (Section~\ref{sec:wakers-ch13}) \\
        \midrule
        \textcolor{orange!70!black}{\textbf{backend}} \newline {\footnotesize\color{gray} chosen at compile time} &
            epoll on Linux, Android and illumos $\cdot$ kqueue on macOS, iOS and the BSDs $\cdot$ event ports on Solaris $\cdot$ IOCP with AFD on Windows $\cdot$ \texttt{poll(2)} on AIX, Fuchsia, Haiku, ESP-IDF, WASI and others \\
        \midrule
        \textbf{kernel} & sockets, pipes, eventfd; the network stack; interrupts \\
        \bottomrule
    \end{tabularx}
    \caption{The layers between an \texttt{.await} and the kernel.}
    \label{fig:wakers-12.1}
\end{figure}

Listing~\ref{lst:wakers-A.1} in Section~\ref{sec:wakers-programs} uses mio on its own, without any async code. Running it under \code{strace} shows every epoll system call it makes:

\begin{lstlisting}[language={}, numbers=none, caption={The epoll system calls made by Listing~\ref{lst:wakers-A.1}. Each token comes back in epoll's \texttt{data} field.}, label=lst:wakers-12.1]
epoll_create1(EPOLL_CLOEXEC) = 3                                      <- Poll::new()
epoll_ctl(3, EPOLL_CTL_ADD, 4, {EPOLLIN|EPOLLRDHUP|EPOLLET, data=0})  <- register listener as Token(0)
epoll_wait(3, [{EPOLLIN, data=0}], 128, -1) = 1                       <- listener is ready
epoll_ctl(3, EPOLL_CTL_ADD, 6, {EPOLLIN|EPOLLRDHUP|EPOLLET, data=1})  <- register socket as Token(1)
epoll_wait(3, [{EPOLLIN, data=1}], 128, -1) = 1                       <- socket is ready
\end{lstlisting}

Two observations:

\begin{itemize}
    \item Tokio's I/O driver is this same loop. Instead of running application code for each token, it follows the token to the socket's state and calls the wakers stored there.
    \item The \code{EPOLLET} flag shows that mio always registers in edge-triggered mode, so whoever reads must keep reading until the operating system says \code{WouldBlock}.
\end{itemize}

\begin{table}[htbp]
    \centering\small
    \begin{tabularx}{\linewidth}{>{\raggedright\arraybackslash}p{4.4cm} Y}
        \toprule
        \textbf{Not mio's job} & \textbf{Who does it} \\
        \midrule
        Reading and writing bytes & The caller; mio only reports when an operation won't block \\
        Timers & Tokio's timer wheel (Section~\ref{sec:wakers-ch15}); mio offers only a timeout argument \\
        Files & No one at this layer (Axiom~\ref{ax:files-always-ready}) \\
        Threads & None are created; \texttt{poll} blocks whichever thread calls it \\
        Futures, tasks and wakers & Tokio, entirely \\
        \bottomrule
    \end{tabularx}
    \caption{The edges of mio's responsibility.}
    \label{tab:wakers-12.1}
\end{table}

\subsubsection{Without epoll}

epoll exists only on Linux. mio picks a backend at compile time, so there is no runtime cost, and Tokio's code above mio is identical everywhere:

\begin{itemize}
    \item \textbf{epoll:} Linux, Android, illumos.
    \item \textbf{kqueue:} macOS, iOS and the BSDs. Same readiness model.
    \item \textbf{Event ports:} Solaris. Same readiness model.
    \item \textbf{IOCP with AFD:} Windows. Readiness built on completion; see below.
    \item \textbf{\code{poll(2)}:} AIX, Fuchsia, Haiku, ESP-IDF, WASI and others. It works, but scans every descriptor on every call.
\end{itemize}

Windows uses the completion model (Definition~\ref{def:completion-model}), but mio promises readiness (Definition~\ref{def:readiness-model}), so on Windows it builds one on top of the other:

\begin{enumerate}
    \item For each socket, mio asks the Winsock kernel driver to complete a request when the socket becomes readable or writable (\code{IOCTL\_AFD\_POLL}, sent to \texttt{\textbackslash Device\textbackslash Afd}).
    \item Those requests finish on an I/O completion port, like any other Windows I/O.
    \item \code{Poll::poll} waits on the port and converts each finished request into a readiness event carrying the original token.
\end{enumerate}

The libuv library beneath Node.js uses the same technique.

\begin{quotebox}{Summary}
    \begin{itemize}
        \item mio answers ``which handles are ready?'' synchronously; Tokio answers ``which task do I wake?''
        \item The token crosses the boundary and the kernel unchanged (Lemma~\ref{lem:token-round-trip}).
        \item mio uses epoll, kqueue, event ports, IOCP with AFD, or \code{poll(2)}, chosen at compile time.
    \end{itemize}
\end{quotebox}

\subsection{Two kinds of wake}
\label{sec:wakers-ch13}

\textit{mio has a type called \code{Waker}. It shares nothing with \code{std::task::Waker} but the name, and understanding both completes the story of step 6 in Section~\ref{sec:wakers-ch5}.}

\begin{definition}{Task-level wake, thread-level wake}{task-thread-level-wake}
    \index{Task-level wake}\index{Thread-level wake}
    A \emph{task-level wake} makes a task runnable by putting it on a run queue; it is what \code{std::task::Waker} does. A \emph{thread-level wake} makes a blocked thread runnable by making a descriptor ready; it is what \code{mio::Waker} does.
\end{definition}

\begin{table}[htbp]
    \centering\small
    \begin{tabularx}{\linewidth}{>{\raggedright\arraybackslash}p{2.4cm} Y Y}
        \toprule
         & \textbf{\texttt{std::task::Waker}} & \textbf{\texttt{mio::Waker}} \\
        \midrule
        Level & Task: ``poll this future again'' & Thread: ``stop sleeping in \texttt{Poll::poll}'' \\
        How many & One per task, cloned freely & One per \texttt{Poll} instance, as mio requires; Tokio registers it as \texttt{TOKEN\_WAKEUP} \\
        What it is & Two pointers (Definition~\ref{def:raw-waker}) & A kernel object registered with the poller \\
        \texttt{wake()} does & Pushes the task onto a run queue & A system call that makes that object ready \\
        \bottomrule
    \end{tabularx}
    \caption{The two wakers.}
    \label{tab:wakers-13.1}
\end{table}

The thread-level object differs by platform:

\begin{itemize}
    \item \textbf{Linux:} an \code{eventfd}; waking writes 1 to it.
    \item \textbf{macOS and the BSDs:} a kqueue user event (\code{EVFILT\_USER}); no descriptor needed.
    \item \textbf{Windows:} a dummy completion posted with \code{PostQueuedCompletionStatus}.
    \item \textbf{\code{poll(2)} platforms:} a pipe.
\end{itemize}

\begin{lemma}{Two-stage remote wake}{two-stage-remote-wake}
    A remote wake that reaches a worker on driver duty consists of a task-level wake followed by a thread-level wake. A local wake needs only the first.
    \why The task-level wake queues the task. By Lemma~\ref{lem:remote-wakes-need-a-poke}, a blocked worker cannot notice that without a descriptor becoming ready, which is what the thread-level wake supplies. A local wake runs on a thread that is already awake.
\end{lemma}

\begin{lstlisting}[language={}, numbers=none, caption={The two stages of a remote wake.}, label=lst:wakers-13.1]
pool thread: job finished
  -> std::task::Waker::wake()          task level: push task B onto the inject queue
    -> the runtime sees that its thread is on driver duty
      -> mio::Waker::wake()            thread level: write 1 to the eventfd
        -> kernel: eventfd readable -> epoll_wait returns with token 0
          -> driver sees TOKEN_WAKEUP; nothing to dispatch, the event only ended the sleep
            -> the thread returns to polling, takes task B, polls it
\end{lstlisting}

\begin{lstlisting}[language={}, numbers=none, caption={\texttt{strace} of a \texttt{current\_thread} runtime awaiting \texttt{spawn\_blocking(|| \{ sleep(50 ms); 42 \})} (Listing~\ref{lst:wakers-A.2}).}, label=lst:wakers-13.2]
eventfd2(0, EFD_CLOEXEC|EFD_NONBLOCK) = 4                    <- mio::Waker::new
epoll_ctl(3, EPOLL_CTL_ADD, 4, {EPOLLIN|EPOLLET, data=0})   <- registered as Token(0)
[main thread]  epoll_wait(3, ..., -1)                       <- driver duty, no timeout
[pool thread]  write(4, "\1\0\0\0\0\0\0\0", 8) = 8          <- job done; task wake -> thread wake
[main thread]  epoll_wait resumed -> [{EPOLLIN, data=0}] = 1 <- woken by token 0, not a socket
\end{lstlisting}

\begin{lemma}{Wake coalescing}{wake-coalescing}
    Any number of thread-level wakes that occur before the sleeping thread runs produce a single event.
    \why Each write only increases the eventfd's counter. The kernel puts a descriptor on epoll's ready list at most once until \code{epoll\_wait} collects it, so the waiting thread receives one event however many writes preceded it.
\end{lemma}

\begin{quotebox}{Summary}
    \begin{itemize}
        \item \code{std::task::Waker} performs task-level wakes; \code{mio::Waker} performs thread-level wakes (Definition~\ref{def:task-thread-level-wake}).
        \item A remote wake to a sleeping worker needs both, in that order (Lemma~\ref{lem:two-stage-remote-wake}). Repeated thread-level wakes coalesce (Lemma~\ref{lem:wake-coalescing}).
    \end{itemize}
\end{quotebox}

\begin{Exercise}[title={Inside Tokio}]
    \begin{enumerate}
        \item Why can't another worker steal the task in a worker's LIFO slot, and what stops two tasks that keep waking each other from monopolising a worker through that slot?
        \item A program builds a \code{current\_thread} runtime, gives its \code{Handle} to another thread that spawns tasks with it, and never calls \code{block\_on}. What happens to those tasks?
        \item In Listing~\ref{lst:wakers-13.2}, why does the eventfd's event carry \code{data=0}? If two pool threads both finish jobs before the sleeping thread runs, how many events does \code{epoll\_wait} return for the eventfd?
        \item In one sentence each, what do \code{Future::poll} and \code{mio::Poll::poll} do, and which of them may block?
    \end{enumerate}
\end{Exercise}
\begin{Answer}
    \begin{enumerate}
        \item The LIFO slot exists so that a task woken by the one just polled runs next, on the same core, while the data it needs is still in cache; only the local run queue is stealable. To stay fair, a worker takes at most three tasks in a row from its LIFO slot before going back to its local queue (Section~\ref{sec:wakers-ch10}).
        \item Nothing runs. A \code{current\_thread} runtime has no worker threads: it polls tasks, handles I/O and fires timers only while some thread is inside its \code{block\_on} (Corollary~\ref{cor:no-block-on-no-progress}). The tasks wait in the inject queue until a thread enters \code{block\_on}.
        \item The eventfd is the \code{mio::Waker}, registered as Token(0), Tokio's \code{TOKEN\_WAKEUP}; the token comes back unchanged in epoll's \code{data} field (Lemma~\ref{lem:token-round-trip}). Two writes before the thread runs produce a single event: each write only adds to the eventfd's counter, and the descriptor sits on the ready list at most once until \code{epoll\_wait} collects it (Lemma~\ref{lem:wake-coalescing}).
        \item \code{Future::poll} asks a future to make as much progress as it can and never blocks; it returns \code{Pending} instead. \code{mio::Poll::poll} is a thin wrapper around \code{epoll\_wait} (or its equivalent) and blocks the calling thread until an event arrives or the timeout passes (Table~\ref{tab:wakers-11.1}).
    \end{enumerate}
\end{Answer}
